\begin{textsample}{NEG Dim 2 – 2016\_c5 – Score -3.00 – 2016\_c5/enade\_2016\_serv\_soc\_b6e03a9959.txt}  \label{ex:f2_neg_032}
CAPÍTULO 5 – DISTRIBUIÇÃO DOS CONCEITOS 5.1 - PANORAMA NACIONAL DA DISTRIBUIÇÃO DOS CONCEITOS O Conceito Enade é um indicador de qualidade que avalia o desempenho dos estudantes a partir dos resultados do Enade.
Os valores possíveis vão de 1 ( pior situação ) a 5 ( melhor situação ).
As Unidades de Observação com apenas um ou sem nenhum concluinte participante não obtêm o Conceito Enade, ficando Sem Conceito ( SC ).
A Tabela 5.1 apresenta a quantidade e distribuição de cursos de Serviço Social participantes do Enade/2016, por faixa de conceito e Grande Região. \textbf{Enfatiza} -se, mais uma vez, que a diferença entre os cursos tabulados neste capítulo e no capítulo 2 corresponde aos cursos sem conceito, ou seja àqueles, em princípio, sem alunos concluintes que houvessem participado da prova ou àqueles que tivessem tão somente um aluno concluinte19.
Observando -se os dados da Tabela 5.1, nota -se que, dos 328 cursos participantes, 138 ( 42,1 \% ) classificaram -se com conceito 3, o valor modal.
Esse foi também o conceito modal em quase todas as regiões : Nordeste ( 40,2 \% ), Sudeste ( 46,0 \% ), Sul ( 39,6 \% ) e Centro-Oeste ( 47,6 \% ).
Na região Norte, a moda foi o conceito 2, atribuído para 46,7 \% dos cursos da região.
O conceito 4 foi o segundo mais frequente em nível nacional ( 24,7 \%, correspondendo a 81 cursos ) e o conceito 2, o terceiro ( 24,4 \%, correspondendo a 80 cursos ).
Houve, ainda, 16 cursos ( 4,9 \% ) que receberam conceito 5 e dez cursos ( 3,0 \% ) que receberam conceito 1.
Dos 328 cursos de Serviço Social, três ( 0,9 \% ) ficaram sem conceito ( SC ). 19 Estes cursos com somente um concluinte não foram considerados no capítulo 2.
A região Norte participou com 30 cursos ou 9,1 \% do total nacional.
Desses, 14 cursos receberam o conceito 2, o que equivale a 46,7 \% do total regional, além de ser o conceito modal na região, como já comentado.
A um curso ( 3,3 \% ) atribuiu -se conceito 1, a nove cursos ( 30,0 \% ), conceito 3 e a cinco cursos ( 16,7 \% ), conceito 4.
Um curso ficou sem conceito ( SC ) e nenhum curso recebeu conceito 5 nessa região.
A região Nordeste participou com 92 cursos ou 28,0 \% do total nacional.
Nessa região, 37 cursos ( 40,2 \% em termos regionais ) obtiveram conceito 3, o conceito modal para a região, como já comentado.
Os conceitos 1 e 2 foram atribuídos, respectivamente, a dois e a 30 cursos ( 2,2 \% e 32,6 \%, respectivamente ).
Ainda nessa região, 15 cursos ( 16,3 \% ) foram avaliados com o conceito 4, e oito cursos ( 8,7 \% ), com conceito 5.
Nenhum curso ficou sem conceito ( SC ).
Dos 137 cursos participantes da região Sudeste ( 41,8 \% do total nacional ), o conceito 3 foi o conceito modal, como já anteriormente assinalado, tendo sido obtido por 63 cursos ( 46,0 \% ).
O conceito 1 foi atribuído a quatro cursos ( 2,9 \% ) e o conceito 2, a 23 cursos ( 16,8 \% ).
Receberam os conceitos 4 e 5, respectivamente, 42 cursos ( 30,7 \% ) e cinco cursos ( 3,6 \% ).
Nessa região, nenhum curso ficou sem conceito ( SC ).
Os 48 cursos da região Sul corresponderam a 14,6 \% do total nacional.
Como já indicado, houve predominância do conceito 3, atribuído a 39,6 \% dos cursos da região ( conceito modal ), ou seja, atribuído a 19 dos 48 cursos participantes na região Sul.
O conceito 1 foi atribuído a dois cursos ( 4,2 \% ), o conceito 2, a sete cursos ( 14,6 \% ), o conceito 4, a 17 cursos ( 35,4 \% ) e o conceito 5, a um curso ( 2,1 \% ).
Dois cursos ficaram sem conceito ( SC ) na região Sul.
Já dos 21 cursos participantes na região Centro-Oeste ( 6,4 \% do total nacional ), dez receberam conceito 3, o conceito modal ( 47,6 \% ), como já fora destacado.
A um curso foi atribuído conceito 1, e a seis cursos ( 28,6 \% ), conceito 2.
Os conceitos restantes, 4 e 5, foram atribuídos a dois cursos, cada.
Nessa região, nenhum curso ficou sem conceito ( SC ).
O Gráfico 5.1 apresenta a distribuição do Conceito Enade segundo Grande Região : poligonais mais à esquerda representam Grandes Regiões com pior distribuição, e poligonais mais à direita, Grandes Regiões com melhores conceitos.
A região Norte ( linha azul ) apresenta os cursos com conceitos mais baixos, ao passo que as regiões Sudeste ( linha vermelha ) e Sul ( linha roxa ) apresentam as distribuições com valores maiores.
Cursos em IES situadas nas regiões Nordeste ( linha verde ) e Centro-Oeste ( linha preta ) apresentam uma situação intermediária. 5.2 - CONCEITOS POR CATEGORIA ADMINISTRATIVA, POR MODALIDADE DE ENSINO E GRANDE REGIÃO A Tabela 5.2 apresenta a distribuição dos cursos participantes do Enade/2016 de Serviço Social, por Categoria Administrativa, de acordo com os conceitos por eles alcançados, segundo as Grandes Regiões brasileiras.
Dos 328 cursos participantes, 62 ( 18,9 \% ) eram ministrados em Instituições Públicas e 266 ( 81,1 \% ), em Instituições Privadas.
De acordo com as informações da Tabela 5.2, em termos nacionais, entre Instituições Públicas e Privadas, observa -se que, dos 16 cursos avaliados com conceito 5, 13 eram oferecidos em IES Públicas e três, em IES Privadas.
Dos 62 cursos participantes de IES Públicas, o conceito 4 foi o valor modal, atribuído a 33 cursos ( 53,2 \% ).
Entre os demais cursos de Instituições Públicas participantes, três cursos foram avaliados com conceito 1, outros três cursos ( 4,8 \% cada ) obtiveram conceito 2, dez cursos ( 16,1 \% ) receberam conceito 3 e, como já comentado, 13 cursos ( 21,0 \% ) receberam conceito 5.
Nesta categoria, nenhum curso ficou sem conceito ( SC ). dez cursos ( 16,1 \% ) receberam conceito 3 e, como já comentado, 13 cursos ( 21,0 \% ) receberam conceito 5.
Nesta categoria, nenhum curso ficou sem conceito ( SC ).
Na rede privada, o conceito modal foi 3, atribuído a 128 cursos, o correspondente a 48,1 \% dos 266 cursos da categoria.
Entre os demais cursos participantes, sete ( 2,6 \% ) receberam conceito 1, e 77 ( 28,9 \% ), conceito 2.
O conceito 4 foi atribuído a 48 cursos ( 18,0 \% ) e o conceito 5, a três cursos ( 1,1 \% ).
Nesta Categoria Administrativa, três cursos ( 1,1 \% ) ficaram sem conceito ( SC ).
Na análise por região, observa -se que, na região Norte, as Instituições Privadas participaram com 23 cursos ( 76,7 \% do total regional ), dos quais um ficou sem conceito ( SC ) e nenhum recebeu conceito 5.
O conceito modal para as Instituições Privadas na região foi 2, atribuído a 13 cursos, correspondendo a 56,5 \% dessa categoria na região.
Dois cursos foram avaliados com conceitos 1 e 4 ( um curso, cada ) e sete cursos, com conceito 3.
As Instituições Públicas participaram com sete cursos na região Norte ( 23,3 \% do total da região ), dos quais quatro obtiveram conceito 4, o valor modal.
O conceito 2 foi atribuído a um curso, e o conceito 3, a dois cursos.
Nenhum curso ficou sem conceito ( SC ) ou com conceito 1 ou 5 nessa combinação de categoria e região.
Na região Nordeste, a rede privada concentrou 74 dos 92 cursos participantes, o equivalente a 80,4 \% do total da região.
Dentre os cursos oferecidos por IES Privadas no Nordeste, 35 cursos foram avaliados com conceito 3, o conceito modal.
Dois cursos obtiveram conceito 1, 30 cursos receberam conceito 2, seis cursos ficaram com conceito 4 e um curso alcançou o conceito máximo, 5.
Nessa combinação de categoria e região, nenhum curso ficou sem conceito ( SC ).
As Instituições Públicas da região Nordeste participaram com 18 cursos ( 19,6 \% do total da região ).
Desses, nove obtiveram conceito 4, o valor modal, seguidos por outros sete cursos, que receberam conceito 5.
O conceito 3 foi atribuído a dois cursos.
Nenhum curso ficou sem conceito ( SC ) ou com conceito 1 ou 2 nessa região.
Na região Sudeste, a proporção de cursos da rede privada, 85,4 \%, foi mais \textbf{elevada} do que a registrada nas demais regiões brasileiras, correspondendo a 117 dos 137 cursos participantes.
Nessa categoria e região, o conceito modal foi 3, atribuído a 61 cursos.
Nenhum curso ficou sem conceito ( SC ).
Os demais foram avaliados com conceito 1 ( três cursos ), conceito 2 ( 22 cursos ), conceito 4 ( 30 cursos ) e conceito 5 ( um curso ).
Entre os 20 cursos oferecidos por Instituições Públicas na região Sudeste, o conceito modal foi 4, atribuído a 12 cursos.
Os demais cursos receberam os conceitos 1 e 2 ( ambos com dois cursos ), 3 ( dois cursos ) e 5 ( quatro cursos ).
Na região Sudeste, nesta categoria, nenhum curso ficou sem conceito ( SC ).
As Instituições Privadas concentraram 35 dos 48 cursos participantes da região Sul, 72,9 \% do total regional, a menor das proporções desta categoria entre as regiões.
Desses, 15 obtiveram conceito 3, o conceito modal.
Nesta combinação de Categoria Administrativa e Grande Região, os conceitos 2 e 4 foram atribuídos a sete e a dez cursos, respectivamente.
Um curso recebeu conceito 5, nenhum recebeu conceito 1 e dois ficaram sem conceito ( SC ).
As Instituições Públicas na região Sul participaram com 13 cursos, 27,1 \% dos cursos da região, a maior das proporções desta categoria nas regiões.
O conceito modal foi 4, atribuído a sete cursos.
Aos demais cursos foram atribuídos os conceitos 3 ( quatro cursos ) e 1 ( dois cursos ).
Nenhum curso da região Sul oferecido por IES Pública ficou sem conceito ( SC ) ou recebeu os conceitos 2 ou 5.
Na região Centro-Oeste, 17 dos 21 cursos participantes eram de Instituições Privadas ( 81,0 \% em termos regionais ).
Desses, dez receberam conceito 3, o conceito modal.
Entre os demais, um recebeu conceito 1, cinco receberam conceito 2 e um recebeu conceito 4.
Nenhum curso ficou sem conceito ( SC ) ou com conceito 5.
Quanto aos quatro cursos oferecidos por Instituições Públicas na região Centro-Oeste ( 19,0 \% do total regional ), dois cursos apresentaram conceito 5, conceito modal, e dois ficaram com conceitos 2 e 4 ( um curso, cada ).
Nenhum curso ficou sem conceito ( SC ) ou com conceitos 1 ou 3.
O Gráfico 5.2 apresenta a distribuição do Conceito Enade segundo a Categoria Administrativa da IES.
Os cursos em IES Públicas ( linha azul ) apresentam uma poligonal mais à direita do que os em IES Privadas ( linha verde ) e, consequentemente, uma distribuição de Conceitos Enade com valores maiores. 5.3 - CONCEITOS POR ORGANIZAÇÃO ACADÊMICA E GRANDE REGIÃO Na Tabela 5.3, encontra -se a distribuição dos conceitos atribuídos aos cursos participantes do Enade/2016 na Área de Serviço Social, por Organização Acadêmica, segundo as Grandes Regiões brasileiras.
Dos 328 cursos de Serviço Social participantes, 145 eram oferecidos em Universidades, 53 em Centros Universitários e 130 em Faculdades.
Esta distribuição corresponde a, respectivamente, 44,2 \%, 16,2 \% e 39,6 \% dos cursos.
De acordo com os dados apresentados, de todos os 13 cursos avaliados com conceito 5, 16 eram vinculados a Universidades.
Os cursos vinculados a esse tipo de Organização Acadêmica tiveram o conceito 3 como conceito modal, atribuído a 50 cursos ( 34,5 \% ).
Os demais cursos vinculados a Universidades receberam os conceitos 1 ( sete cursos ), 2 ( 25 cursos ), 4 ( 48 cursos ) e conceito 5 ( 13 cursos, como já mencionado ).
Ficaram sem conceito ( SC ), dois cursos.
Entre os cursos em Centros Universitários, o conceito modal foi também 3, atribuído a 24 cursos ( 45,3 \% ).
Nenhum curso vinculado a esse tipo de Organização Acadêmica ficou sem conceito ( SC ), e os outros cursos receberam os conceitos 1 ( dois cursos ), 2 ( nove cursos ), 4 ( 15 cursos ) e 5 ( três cursos ).
Dos 130 cursos mantidos por Faculdades, 64 ( 49,2 \% ) receberam conceito 3, o conceito modal.
O segundo conceito mais frequentemente atribuído aos cursos oferecidos por essa categoria de IES foi o conceito 2, recebido por 46 cursos.
Na sequência, 18 cursos obtiveram conceito 4 e um curso, conceito 1.
Nenhum curso foi atribuído ao conceito 5 e um curso ficou sem conceito ( SC ).
Considerando -se separadamente as regiões brasileiras, verifica -se que, na região Norte, as Universidades participaram com 10 dos 30 cursos participantes.
Quatro cursos receberam conceito 4, conceito modal.
Os demais cursos receberam os conceitos 1 ( um curso ), 2 ( dois cursos ) e 3 ( três cursos ).
Nenhum curso ficou sem conceito ( SC ) ou recebeu conceito 5 nesta combinação de categoria e região.
Os Centros Universitários da região Norte foram representados por cinco cursos, aos quais foram atribuídos os conceitos 2 ( valor modal, recebido por três cursos ) e 3 ( recebido por dois cursos ).
Nenhum curso ficou sem conceito ( SC ) ou recebeu os demais conceitos nesta combinação de categoria e região.
As Faculdades participaram com 15 cursos na região Norte que receberam os conceitos 2 ( nove cursos, conceito modal ), 3 ( quatro cursos ) e 4 ( um curso ).
Um curso ficou sem conceito ( SC ).
Nenhuma Faculdade da região obteve conceito 1 ou 5.
Na região Nordeste, as Universidades participaram com 30 dos 92 cursos da Área de Serviço Social oferecidos na região, dos quais nenhum ficou sem conceito ( SC ).
Nota -se que tais cursos apresentaram uma distribuição com moda \textbf{distribuída} em dois conceitos contíguos, 3 e 4, atribuídos, cada qual, a oito cursos.
Obtiveram os conceitos 2 e 5, respectivamente, seis e sete cursos.
Nenhum curso oferecido por Universidades no Nordeste ficou sem conceito ( SC ) e um curso recebeu conceito 1.
Os Centros Universitários contaram com nove cursos participante na região Nordeste.
Esses obtiveram conceitos 2, 4 e 5 ( um curso, cada ) e 3 ( seis cursos, o conceito modal ).
Nenhum curso oferecido por Centros Universitários no Nordeste ficou sem conceito ( SC ) ou recebeu conceito 1.
As Faculdades foram representadas por 53 cursos na região Nordeste, nota -se que tais cursos, também, apresentaram distribuição bimodal, uma vez que as maiores frequências se concentraram em dois conceitos, os conceitos 2 e 3, atribuídos, cada qual, a 23 cursos.
O conceito 1 foi atribuído a um curso, e o conceito 4, a seis cursos.
Na região Nordeste, nenhum curso oferecido por Faculdades recebeu conceito 5, tampouco ficou sem conceito ( SC ).
Na região Sudeste, as Universidades concentraram 65 dos 137 cursos de Serviço Social da região.
Entre os cursos em Universidades, nessa região, o conceito modal foi 3, atribuído a 26 cursos.
Os demais cursos receberam os conceitos 1 ( dois cursos ), 2 ( 12 cursos ), 4 ( 21 cursos ) e 5 ( quatro cursos ).
Nenhum curso ficou sem conceito ( SC ).
Os Centros Universitários participaram com 31 cursos na região Sudeste, dos quais 13 obtiveram o conceito modal 3.
O conceito 1 foi atribuído a dois cursos, o conceito 2, a quatro cursos, o conceito 4, a 11 cursos e, por fim, o conceito 5, a um curso.
Nenhum curso ficou sem conceito ( SC ) nessa combinação de organização acadêmica e região.
Já as Faculdades foram representadas por 41 cursos na região Sudeste, dos quais nenhum ficou sem conceito ( SC ) ou recebeu conceito 1 ou 5.
Aos demais cursos foram atribuídos os conceitos 2 ( sete cursos ), 3 ( 24 cursos, o conceito modal ) e 4 ( dez cursos ).
Dos 48 cursos da região Sul, 29 eram de Universidades, para os quais o conceito modal foi 4, atribuído a 14 cursos.
Nesse tipo de organização, dois cursos ficaram sem conceito ( SC ) e nenhum curso recebeu conceito 5.
Os demais cursos receberam os conceitos 1 e 2 ( dois cursos, cada ) e 3 ( nove cursos ).
Dos cinco cursos participantes de Centros Universitários da região Sul, nenhum ficou sem conceito ( SC ) ou recebeu conceito 1.
O conceito modal, 4, foi atribuído a dois cursos.
Os conceitos restantes, 2, 3 e 5, foram atribuídos a um curso, cada.
Quanto aos 14 cursos vinculados a Faculdades na região Sul, nove receberam o conceito modal 3.
O conceito 2 foi atribuído a quatro cursos, e o conceito 4, a um curso.
Nenhum curso ficou sem conceito ( SC ) ou recebeu conceito 1ou 5.
Na região Centro-Oeste, 11 dos 21 cursos eram de Universidades.
Nesse tipo de organização, nenhum curso ficou sem conceito ( SC ).
Os conceitos 1 e 4 foram atribuídos a um curso cada, o conceito 2, a três cursos e o conceito 3, conceito modal, a quatro cursos.
Além disso, dois cursos alcançaram o conceito máximo, 5.
Os Centros Universitários da região Centro-Oeste contaram com três cursos, que obtiveram os conceitos 3 ( dois cursos, conceito modal ) e 4 ( um curso ).
Dos sete cursos oferecidos em Faculdades na região Centro-Oeste, quatro receberam o conceito modal 3, e três cursos ficaram com conceito 2.
O Gráfico 5.3 apresenta a distribuição do Conceito Enade segundo a Organização Acadêmica da IES.
Os cursos em Universidades ( linha azul ) apresentam uma poligonal mais à direita do que as demais, denotando uma melhor distribuição de Conceitos Enade.
No outro extremo, temos os cursos oferecidos em Faculdades ( linha vermelha ) com a poligonal mais à esquerda e, consequentemente, uma distribuição de Conceitos Enade com menores valores.
IES de Centros Universitários ( linha verde ) apresentam uma situação intermediária.

% matched lemmas: distribuir, elevado, enfatizar
\end{textsample}
